% =====================================================================
% Figura: apantallamiento electrostático (cavidad en un conductor).
% (a) Corteza conductora neutra (radios a<b) con una carga q excéntrica en
%     la cavidad y una carga exterior q_e. Líneas de campo exactas (imágenes):
%     dentro, q + imagen respecto de r=a; fuera, q_e + imagen respecto de
%     r=b + carga central (q - q_1). Las marcas -/+ se colocan donde nacen o
%     mueren las líneas: su densidad es proporcional a |sigma|.
% (b) Cavidad vacía en un conductor a tierra: Dirichlet + Laplace => phi=0.
% =====================================================================
\begin{tikzpicture}
  \def\esc{1.2}
  % ---------------- (a)
  \begin{scope}
    \clip (-4.3*\esc,-2.45*\esc) rectangle (3.4*\esc,2.45*\esc);
    \fill[panel] (-4.3*\esc,-2.45*\esc) rectangle (3.4*\esc,2.45*\esc);
    % corteza conductora (anillo) con sombreado radial sobrio
    \fill[even odd rule,inner color=conductor,outer color=conductorosc]
      (0,0) circle (1.5*\esc) (0,0) circle (1.0*\esc);
    \draw[gristexto,line width=0.5pt] (0,0) circle (1.5*\esc) (0,0) circle (1.0*\esc);
    \fill[fondo] (0,0) circle (1.0*\esc);
    % Generado automáticamente por generar_figuras.py; no editar.
\draw[flinea, flecha=0.352cm] plot coordinates {(0.518,0.037) (0.559,0.062) (0.600,0.086) (0.642,0.110) (0.683,0.134) (0.725,0.158) (0.767,0.181) (0.810,0.203) (0.852,0.225) (0.895,0.246) (0.939,0.266) (0.983,0.286) (1.028,0.304) (1.073,0.320) (1.118,0.336) (1.149,0.346)};
\draw[flinea, flecha=0.352cm] plot coordinates {(0.518,-0.037) (0.559,-0.062) (0.600,-0.086) (0.642,-0.110) (0.683,-0.134) (0.725,-0.158) (0.767,-0.181) (0.810,-0.203) (0.852,-0.225) (0.895,-0.246) (0.939,-0.266) (0.983,-0.286) (1.028,-0.304) (1.073,-0.320) (1.118,-0.336) (1.149,-0.346)};
\draw[flinea, flecha=0.389cm] plot coordinates {(0.497,0.059) (0.525,0.098) (0.553,0.138) (0.581,0.176) (0.609,0.215) (0.638,0.253) (0.668,0.291) (0.698,0.328) (0.729,0.365) (0.762,0.400) (0.795,0.435) (0.829,0.469) (0.864,0.501) (0.901,0.532) (0.938,0.562) (0.977,0.590) (1.018,0.616) (1.026,0.622)};
\draw[flinea, flecha=0.389cm] plot coordinates {(0.497,-0.059) (0.525,-0.098) (0.553,-0.138) (0.581,-0.176) (0.609,-0.215) (0.638,-0.253) (0.668,-0.291) (0.698,-0.328) (0.729,-0.365) (0.762,-0.400) (0.795,-0.435) (0.829,-0.469) (0.864,-0.501) (0.901,-0.532) (0.938,-0.562) (0.977,-0.590) (1.018,-0.616) (1.026,-0.622)};
\draw[flinea, flecha=0.433cm] plot coordinates {(0.477,0.069) (0.490,0.115) (0.505,0.161) (0.519,0.207) (0.534,0.252) (0.550,0.297) (0.567,0.342) (0.584,0.387) (0.603,0.432) (0.622,0.475) (0.643,0.519) (0.665,0.561) (0.689,0.603) (0.714,0.644) (0.741,0.684) (0.769,0.722) (0.799,0.760) (0.831,0.796) (0.865,0.830) (0.866,0.831)};
\draw[flinea, flecha=0.433cm] plot coordinates {(0.477,-0.069) (0.490,-0.115) (0.505,-0.161) (0.519,-0.207) (0.534,-0.252) (0.550,-0.297) (0.567,-0.342) (0.584,-0.387) (0.603,-0.432) (0.622,-0.475) (0.643,-0.519) (0.665,-0.561) (0.689,-0.603) (0.714,-0.644) (0.741,-0.684) (0.769,-0.722) (0.799,-0.760) (0.831,-0.796) (0.865,-0.830) (0.866,-0.831)};
\draw[flinea, flecha=0.485cm] plot coordinates {(0.456,0.072) (0.456,0.120) (0.457,0.168) (0.458,0.216) (0.459,0.264) (0.461,0.312) (0.464,0.360) (0.468,0.408) (0.472,0.455) (0.478,0.503) (0.485,0.551) (0.494,0.598) (0.503,0.645) (0.515,0.691) (0.528,0.738) (0.542,0.783) (0.559,0.828) (0.577,0.873) (0.598,0.916) (0.620,0.959) (0.645,1.000) (0.650,1.009)};
\draw[flinea, flecha=0.485cm] plot coordinates {(0.456,-0.072) (0.456,-0.120) (0.457,-0.168) (0.458,-0.216) (0.459,-0.264) (0.461,-0.312) (0.464,-0.360) (0.468,-0.408) (0.472,-0.455) (0.478,-0.503) (0.485,-0.551) (0.494,-0.598) (0.503,-0.645) (0.515,-0.691) (0.528,-0.738) (0.542,-0.783) (0.559,-0.828) (0.577,-0.873) (0.598,-0.916) (0.620,-0.959) (0.645,-1.000) (0.650,-1.009)};
\draw[flinea, flecha=0.549cm] plot coordinates {(0.435,0.069) (0.422,0.115) (0.409,0.161) (0.396,0.207) (0.383,0.254) (0.371,0.300) (0.360,0.347) (0.350,0.394) (0.340,0.441) (0.331,0.488) (0.323,0.535) (0.316,0.583) (0.311,0.630) (0.307,0.678) (0.304,0.726) (0.302,0.774) (0.302,0.822) (0.304,0.870) (0.307,0.918) (0.313,0.966) (0.320,1.013) (0.329,1.060) (0.340,1.107) (0.351,1.147)};
\draw[flinea, flecha=0.549cm] plot coordinates {(0.435,-0.069) (0.422,-0.115) (0.409,-0.161) (0.396,-0.207) (0.383,-0.254) (0.371,-0.300) (0.360,-0.347) (0.350,-0.394) (0.340,-0.441) (0.331,-0.488) (0.323,-0.535) (0.316,-0.583) (0.311,-0.630) (0.307,-0.678) (0.304,-0.726) (0.302,-0.774) (0.302,-0.822) (0.304,-0.870) (0.307,-0.918) (0.313,-0.966) (0.320,-1.013) (0.329,-1.060) (0.340,-1.107) (0.351,-1.147)};
\draw[flinea, flecha=0.629cm] plot coordinates {(0.415,0.059) (0.388,0.099) (0.360,0.138) (0.334,0.178) (0.307,0.218) (0.281,0.258) (0.255,0.299) (0.230,0.339) (0.205,0.381) (0.181,0.422) (0.157,0.464) (0.135,0.506) (0.113,0.549) (0.091,0.592) (0.071,0.635) (0.052,0.679) (0.034,0.724) (0.016,0.769) (0.000,0.814) (-0.015,0.859) (-0.029,0.905) (-0.041,0.952) (-0.052,0.998) (-0.061,1.046) (-0.069,1.093) (-0.075,1.141) (-0.080,1.188) (-0.080,1.197)};
\draw[flinea, flecha=0.629cm] plot coordinates {(0.415,-0.059) (0.388,-0.099) (0.360,-0.138) (0.334,-0.178) (0.307,-0.218) (0.281,-0.258) (0.255,-0.299) (0.230,-0.339) (0.205,-0.381) (0.181,-0.422) (0.157,-0.464) (0.135,-0.506) (0.113,-0.549) (0.091,-0.592) (0.071,-0.635) (0.052,-0.679) (0.034,-0.724) (0.016,-0.769) (0.000,-0.814) (-0.015,-0.859) (-0.029,-0.905) (-0.041,-0.952) (-0.052,-0.998) (-0.061,-1.046) (-0.069,-1.093) (-0.075,-1.141) (-0.080,-1.188) (-0.080,-1.197)};
\draw[flinea, flecha=0.730cm] plot coordinates {(0.394,0.037) (0.353,0.062) (0.312,0.087) (0.271,0.112) (0.230,0.137) (0.190,0.163) (0.149,0.188) (0.109,0.214) (0.069,0.241) (0.029,0.267) (-0.011,0.294) (-0.050,0.321) (-0.090,0.349) (-0.129,0.377) (-0.167,0.406) (-0.205,0.435) (-0.243,0.464) (-0.281,0.494) (-0.318,0.524) (-0.355,0.555) (-0.391,0.587) (-0.427,0.619) (-0.462,0.651) (-0.497,0.684) (-0.532,0.718) (-0.565,0.752) (-0.598,0.787) (-0.631,0.822) (-0.662,0.858) (-0.693,0.895) (-0.724,0.932) (-0.736,0.948)};
\draw[flinea, flecha=0.730cm] plot coordinates {(0.394,-0.037) (0.353,-0.062) (0.312,-0.087) (0.271,-0.112) (0.230,-0.137) (0.190,-0.163) (0.149,-0.188) (0.109,-0.214) (0.069,-0.241) (0.029,-0.267) (-0.011,-0.294) (-0.050,-0.321) (-0.090,-0.349) (-0.129,-0.377) (-0.167,-0.406) (-0.205,-0.435) (-0.243,-0.464) (-0.281,-0.494) (-0.318,-0.524) (-0.355,-0.555) (-0.391,-0.587) (-0.427,-0.619) (-0.462,-0.651) (-0.497,-0.684) (-0.532,-0.718) (-0.565,-0.752) (-0.598,-0.787) (-0.631,-0.822) (-0.662,-0.858) (-0.693,-0.895) (-0.724,-0.932) (-0.736,-0.948)};
\draw[flinea, flecha=1.728cm] plot coordinates {(-3.884,0.037) (-3.841,0.058) (-3.798,0.078) (-3.755,0.099) (-3.712,0.120) (-3.668,0.141) (-3.625,0.162) (-3.582,0.183) (-3.539,0.204) (-3.496,0.226) (-3.453,0.247) (-3.410,0.268) (-3.367,0.289) (-3.324,0.310) (-3.281,0.332) (-3.238,0.353) (-3.195,0.375) (-3.152,0.396) (-3.109,0.418) (-3.066,0.439) (-3.023,0.461) (-2.980,0.483) (-2.938,0.504) (-2.895,0.526) (-2.852,0.548) (-2.810,0.570) (-2.767,0.592) (-2.724,0.614) (-2.682,0.636) (-2.639,0.658) (-2.596,0.680) (-2.554,0.703) (-2.511,0.725) (-2.469,0.747) (-2.426,0.769) (-2.383,0.791) (-2.340,0.812) (-2.298,0.834) (-2.255,0.856) (-2.212,0.877) (-2.169,0.898) (-2.126,0.919) (-2.082,0.940) (-2.039,0.961) (-1.996,0.981) (-1.952,1.002) (-1.909,1.022) (-1.865,1.043) (-1.822,1.063) (-1.779,1.085) (-1.737,1.108) (-1.696,1.133) (-1.657,1.161) (-1.622,1.194) (-1.594,1.232) (-1.572,1.275) (-1.557,1.321) (-1.546,1.368) (-1.539,1.415) (-1.533,1.463) (-1.528,1.510) (-1.525,1.558) (-1.522,1.606) (-1.519,1.654) (-1.516,1.702) (-1.513,1.750) (-1.511,1.798) (-1.508,1.846) (-1.506,1.894) (-1.503,1.942) (-1.500,1.990) (-1.497,2.037) (-1.494,2.085) (-1.491,2.133) (-1.487,2.181) (-1.484,2.229) (-1.481,2.277) (-1.477,2.325) (-1.473,2.373) (-1.469,2.420) (-1.465,2.468) (-1.461,2.516) (-1.457,2.564) (-1.453,2.612) (-1.449,2.660) (-1.444,2.707) (-1.440,2.755) (-1.435,2.803) (-1.430,2.851) (-1.426,2.898) (-1.421,2.946)};
\draw[flinea, flecha=1.728cm] plot coordinates {(-3.884,-0.037) (-3.841,-0.058) (-3.798,-0.078) (-3.755,-0.099) (-3.712,-0.120) (-3.668,-0.141) (-3.625,-0.162) (-3.582,-0.183) (-3.539,-0.204) (-3.496,-0.226) (-3.453,-0.247) (-3.410,-0.268) (-3.367,-0.289) (-3.324,-0.310) (-3.281,-0.332) (-3.238,-0.353) (-3.195,-0.375) (-3.152,-0.396) (-3.109,-0.418) (-3.066,-0.439) (-3.023,-0.461) (-2.980,-0.483) (-2.938,-0.504) (-2.895,-0.526) (-2.852,-0.548) (-2.810,-0.570) (-2.767,-0.592) (-2.724,-0.614) (-2.682,-0.636) (-2.639,-0.658) (-2.596,-0.680) (-2.554,-0.703) (-2.511,-0.725) (-2.469,-0.747) (-2.426,-0.769) (-2.383,-0.791) (-2.340,-0.812) (-2.298,-0.834) (-2.255,-0.856) (-2.212,-0.877) (-2.169,-0.898) (-2.126,-0.919) (-2.082,-0.940) (-2.039,-0.961) (-1.996,-0.981) (-1.952,-1.002) (-1.909,-1.022) (-1.865,-1.043) (-1.822,-1.063) (-1.779,-1.085) (-1.737,-1.108) (-1.696,-1.133) (-1.657,-1.161) (-1.622,-1.194) (-1.594,-1.232) (-1.572,-1.275) (-1.557,-1.321) (-1.546,-1.368) (-1.539,-1.415) (-1.533,-1.463) (-1.528,-1.510) (-1.525,-1.558) (-1.522,-1.606) (-1.519,-1.654) (-1.516,-1.702) (-1.513,-1.750) (-1.511,-1.798) (-1.508,-1.846) (-1.506,-1.894) (-1.503,-1.942) (-1.500,-1.990) (-1.497,-2.037) (-1.494,-2.085) (-1.491,-2.133) (-1.487,-2.181) (-1.484,-2.229) (-1.481,-2.277) (-1.477,-2.325) (-1.473,-2.373) (-1.469,-2.420) (-1.465,-2.468) (-1.461,-2.516) (-1.457,-2.564) (-1.453,-2.612) (-1.449,-2.660) (-1.444,-2.707) (-1.440,-2.755) (-1.435,-2.803) (-1.430,-2.851) (-1.426,-2.898) (-1.421,-2.946)};
\draw[flinea, flecha=1.382cm] plot coordinates {(-3.901,0.060) (-3.868,0.094) (-3.834,0.129) (-3.800,0.163) (-3.767,0.197) (-3.733,0.232) (-3.700,0.266) (-3.666,0.300) (-3.633,0.335) (-3.600,0.369) (-3.566,0.404) (-3.533,0.438) (-3.500,0.473) (-3.467,0.508) (-3.434,0.543) (-3.401,0.577) (-3.368,0.612) (-3.335,0.647) (-3.302,0.683) (-3.270,0.718) (-3.237,0.753) (-3.205,0.789) (-3.172,0.824) (-3.140,0.860) (-3.108,0.896) (-3.077,0.932) (-3.045,0.968) (-3.014,1.004) (-2.982,1.041) (-2.952,1.077) (-2.921,1.114) (-2.890,1.151) (-2.860,1.189) (-2.831,1.226) (-2.801,1.264) (-2.772,1.302) (-2.743,1.341) (-2.715,1.380) (-2.687,1.419) (-2.660,1.459) (-2.634,1.498) (-2.608,1.539) (-2.582,1.580) (-2.558,1.621) (-2.534,1.662) (-2.510,1.704) (-2.488,1.747) (-2.466,1.789) (-2.445,1.833) (-2.425,1.876) (-2.406,1.920) (-2.387,1.965) (-2.370,2.009) (-2.353,2.054) (-2.337,2.099) (-2.321,2.145) (-2.307,2.191) (-2.293,2.237) (-2.280,2.283) (-2.267,2.329) (-2.255,2.375) (-2.244,2.422) (-2.233,2.469) (-2.223,2.516) (-2.213,2.563) (-2.203,2.610) (-2.194,2.657) (-2.186,2.704) (-2.178,2.751) (-2.170,2.799) (-2.162,2.846) (-2.155,2.894) (-2.148,2.941)};
\draw[flinea, flecha=1.382cm] plot coordinates {(-3.901,-0.060) (-3.868,-0.094) (-3.834,-0.129) (-3.800,-0.163) (-3.767,-0.197) (-3.733,-0.232) (-3.700,-0.266) (-3.666,-0.300) (-3.633,-0.335) (-3.600,-0.369) (-3.566,-0.404) (-3.533,-0.438) (-3.500,-0.473) (-3.467,-0.508) (-3.434,-0.543) (-3.401,-0.577) (-3.368,-0.612) (-3.335,-0.647) (-3.302,-0.683) (-3.270,-0.718) (-3.237,-0.753) (-3.205,-0.789) (-3.172,-0.824) (-3.140,-0.860) (-3.108,-0.896) (-3.077,-0.932) (-3.045,-0.968) (-3.014,-1.004) (-2.982,-1.041) (-2.952,-1.077) (-2.921,-1.114) (-2.890,-1.151) (-2.860,-1.189) (-2.831,-1.226) (-2.801,-1.264) (-2.772,-1.302) (-2.743,-1.341) (-2.715,-1.380) (-2.687,-1.419) (-2.660,-1.459) (-2.634,-1.498) (-2.608,-1.539) (-2.582,-1.580) (-2.558,-1.621) (-2.534,-1.662) (-2.510,-1.704) (-2.488,-1.747) (-2.466,-1.789) (-2.445,-1.833) (-2.425,-1.876) (-2.406,-1.920) (-2.387,-1.965) (-2.370,-2.009) (-2.353,-2.054) (-2.337,-2.099) (-2.321,-2.145) (-2.307,-2.191) (-2.293,-2.237) (-2.280,-2.283) (-2.267,-2.329) (-2.255,-2.375) (-2.244,-2.422) (-2.233,-2.469) (-2.223,-2.516) (-2.213,-2.563) (-2.203,-2.610) (-2.194,-2.657) (-2.186,-2.704) (-2.178,-2.751) (-2.170,-2.799) (-2.162,-2.846) (-2.155,-2.894) (-2.148,-2.941)};
\draw[flinea, flecha=1.243cm] plot coordinates {(-3.918,0.073) (-3.894,0.114) (-3.870,0.156) (-3.846,0.198) (-3.822,0.239) (-3.798,0.281) (-3.774,0.322) (-3.751,0.364) (-3.727,0.406) (-3.703,0.448) (-3.680,0.489) (-3.656,0.531) (-3.633,0.573) (-3.609,0.615) (-3.586,0.657) (-3.563,0.699) (-3.540,0.741) (-3.517,0.783) (-3.494,0.826) (-3.472,0.868) (-3.449,0.910) (-3.427,0.953) (-3.405,0.996) (-3.383,1.038) (-3.361,1.081) (-3.340,1.124) (-3.319,1.167) (-3.298,1.210) (-3.277,1.254) (-3.257,1.297) (-3.237,1.341) (-3.217,1.384) (-3.197,1.428) (-3.178,1.472) (-3.160,1.517) (-3.141,1.561) (-3.123,1.605) (-3.106,1.650) (-3.089,1.695) (-3.072,1.740) (-3.056,1.785) (-3.040,1.830) (-3.025,1.876) (-3.010,1.922) (-2.996,1.967) (-2.982,2.013) (-2.968,2.060) (-2.956,2.106) (-2.943,2.152) (-2.932,2.199) (-2.920,2.245) (-2.909,2.292) (-2.899,2.339) (-2.889,2.386) (-2.880,2.433) (-2.871,2.480) (-2.862,2.527) (-2.854,2.575) (-2.846,2.622) (-2.839,2.670) (-2.832,2.717) (-2.826,2.765) (-2.819,2.812) (-2.814,2.860) (-2.808,2.908) (-2.804,2.943)};
\draw[flinea, flecha=1.243cm] plot coordinates {(-3.918,-0.073) (-3.894,-0.114) (-3.870,-0.156) (-3.846,-0.198) (-3.822,-0.239) (-3.798,-0.281) (-3.774,-0.322) (-3.751,-0.364) (-3.727,-0.406) (-3.703,-0.448) (-3.680,-0.489) (-3.656,-0.531) (-3.633,-0.573) (-3.609,-0.615) (-3.586,-0.657) (-3.563,-0.699) (-3.540,-0.741) (-3.517,-0.783) (-3.494,-0.826) (-3.472,-0.868) (-3.449,-0.910) (-3.427,-0.953) (-3.405,-0.996) (-3.383,-1.038) (-3.361,-1.081) (-3.340,-1.124) (-3.319,-1.167) (-3.298,-1.210) (-3.277,-1.254) (-3.257,-1.297) (-3.237,-1.341) (-3.217,-1.384) (-3.197,-1.428) (-3.178,-1.472) (-3.160,-1.517) (-3.141,-1.561) (-3.123,-1.605) (-3.106,-1.650) (-3.089,-1.695) (-3.072,-1.740) (-3.056,-1.785) (-3.040,-1.830) (-3.025,-1.876) (-3.010,-1.922) (-2.996,-1.967) (-2.982,-2.013) (-2.968,-2.060) (-2.956,-2.106) (-2.943,-2.152) (-2.932,-2.199) (-2.920,-2.245) (-2.909,-2.292) (-2.899,-2.339) (-2.889,-2.386) (-2.880,-2.433) (-2.871,-2.480) (-2.862,-2.527) (-2.854,-2.575) (-2.846,-2.622) (-2.839,-2.670) (-2.832,-2.717) (-2.826,-2.765) (-2.819,-2.812) (-2.814,-2.860) (-2.808,-2.908) (-2.804,-2.943)};
\draw[flinea, flecha=1.171cm] plot coordinates {(-3.935,0.080) (-3.920,0.126) (-3.906,0.172) (-3.892,0.218) (-3.877,0.263) (-3.863,0.309) (-3.849,0.355) (-3.835,0.401) (-3.821,0.447) (-3.807,0.493) (-3.793,0.539) (-3.779,0.585) (-3.765,0.631) (-3.751,0.677) (-3.738,0.723) (-3.725,0.769) (-3.711,0.815) (-3.698,0.861) (-3.685,0.907) (-3.672,0.953) (-3.660,1.000) (-3.647,1.046) (-3.635,1.093) (-3.623,1.139) (-3.611,1.186) (-3.599,1.232) (-3.588,1.279) (-3.577,1.325) (-3.566,1.372) (-3.555,1.419) (-3.545,1.466) (-3.535,1.513) (-3.525,1.560) (-3.516,1.607) (-3.507,1.654) (-3.498,1.701) (-3.489,1.748) (-3.481,1.796) (-3.473,1.843) (-3.466,1.890) (-3.459,1.938) (-3.452,1.985) (-3.445,2.033) (-3.439,2.081) (-3.434,2.128) (-3.428,2.176) (-3.423,2.224) (-3.418,2.271) (-3.414,2.319) (-3.410,2.367) (-3.406,2.415) (-3.403,2.463) (-3.400,2.511) (-3.397,2.559) (-3.395,2.607) (-3.393,2.655) (-3.391,2.703) (-3.390,2.751) (-3.389,2.799) (-3.388,2.847) (-3.387,2.895) (-3.387,2.943)};
\draw[flinea, flecha=1.171cm] plot coordinates {(-3.935,-0.080) (-3.920,-0.126) (-3.906,-0.172) (-3.892,-0.218) (-3.877,-0.263) (-3.863,-0.309) (-3.849,-0.355) (-3.835,-0.401) (-3.821,-0.447) (-3.807,-0.493) (-3.793,-0.539) (-3.779,-0.585) (-3.765,-0.631) (-3.751,-0.677) (-3.738,-0.723) (-3.725,-0.769) (-3.711,-0.815) (-3.698,-0.861) (-3.685,-0.907) (-3.672,-0.953) (-3.660,-1.000) (-3.647,-1.046) (-3.635,-1.093) (-3.623,-1.139) (-3.611,-1.186) (-3.599,-1.232) (-3.588,-1.279) (-3.577,-1.325) (-3.566,-1.372) (-3.555,-1.419) (-3.545,-1.466) (-3.535,-1.513) (-3.525,-1.560) (-3.516,-1.607) (-3.507,-1.654) (-3.498,-1.701) (-3.489,-1.748) (-3.481,-1.796) (-3.473,-1.843) (-3.466,-1.890) (-3.459,-1.938) (-3.452,-1.985) (-3.445,-2.033) (-3.439,-2.081) (-3.434,-2.128) (-3.428,-2.176) (-3.423,-2.224) (-3.418,-2.271) (-3.414,-2.319) (-3.410,-2.367) (-3.406,-2.415) (-3.403,-2.463) (-3.400,-2.511) (-3.397,-2.559) (-3.395,-2.607) (-3.393,-2.655) (-3.391,-2.703) (-3.390,-2.751) (-3.389,-2.799) (-3.388,-2.847) (-3.387,-2.895) (-3.387,-2.943)};
\draw[flinea, flecha=1.147cm] plot coordinates {(-3.952,0.084) (-3.947,0.131) (-3.942,0.179) (-3.937,0.227) (-3.933,0.275) (-3.928,0.322) (-3.923,0.370) (-3.919,0.418) (-3.914,0.466) (-3.910,0.514) (-3.906,0.561) (-3.901,0.609) (-3.897,0.657) (-3.893,0.705) (-3.889,0.753) (-3.886,0.801) (-3.882,0.848) (-3.878,0.896) (-3.875,0.944) (-3.872,0.992) (-3.869,1.040) (-3.866,1.088) (-3.863,1.136) (-3.861,1.184) (-3.858,1.232) (-3.856,1.280) (-3.854,1.328) (-3.853,1.376) (-3.851,1.424) (-3.850,1.472) (-3.849,1.519) (-3.848,1.567) (-3.848,1.615) (-3.847,1.663) (-3.847,1.711) (-3.848,1.759) (-3.848,1.807) (-3.849,1.855) (-3.850,1.903) (-3.851,1.951) (-3.852,1.999) (-3.854,2.047) (-3.856,2.095) (-3.859,2.143) (-3.861,2.191) (-3.864,2.239) (-3.867,2.287) (-3.870,2.335) (-3.874,2.383) (-3.878,2.431) (-3.882,2.478) (-3.886,2.526) (-3.891,2.574) (-3.895,2.622) (-3.900,2.670) (-3.906,2.717) (-3.911,2.765) (-3.917,2.813) (-3.923,2.860) (-3.929,2.908) (-3.933,2.944)};
\draw[flinea, flecha=1.147cm] plot coordinates {(-3.952,-0.084) (-3.947,-0.131) (-3.942,-0.179) (-3.937,-0.227) (-3.933,-0.275) (-3.928,-0.322) (-3.923,-0.370) (-3.919,-0.418) (-3.914,-0.466) (-3.910,-0.514) (-3.906,-0.561) (-3.901,-0.609) (-3.897,-0.657) (-3.893,-0.705) (-3.889,-0.753) (-3.886,-0.801) (-3.882,-0.848) (-3.878,-0.896) (-3.875,-0.944) (-3.872,-0.992) (-3.869,-1.040) (-3.866,-1.088) (-3.863,-1.136) (-3.861,-1.184) (-3.858,-1.232) (-3.856,-1.280) (-3.854,-1.328) (-3.853,-1.376) (-3.851,-1.424) (-3.850,-1.472) (-3.849,-1.519) (-3.848,-1.567) (-3.848,-1.615) (-3.847,-1.663) (-3.847,-1.711) (-3.848,-1.759) (-3.848,-1.807) (-3.849,-1.855) (-3.850,-1.903) (-3.851,-1.951) (-3.852,-1.999) (-3.854,-2.047) (-3.856,-2.095) (-3.859,-2.143) (-3.861,-2.191) (-3.864,-2.239) (-3.867,-2.287) (-3.870,-2.335) (-3.874,-2.383) (-3.878,-2.431) (-3.882,-2.478) (-3.886,-2.526) (-3.891,-2.574) (-3.895,-2.622) (-3.900,-2.670) (-3.906,-2.717) (-3.911,-2.765) (-3.917,-2.813) (-3.923,-2.860) (-3.929,-2.908) (-3.933,-2.944)};
\draw[flinea, flecha=1.166cm] plot coordinates {(-3.968,0.084) (-3.973,0.131) (-3.978,0.179) (-3.983,0.227) (-3.988,0.275) (-3.993,0.322) (-3.998,0.370) (-4.003,0.418) (-4.008,0.466) (-4.013,0.513) (-4.018,0.561) (-4.024,0.609) (-4.029,0.656) (-4.035,0.704) (-4.040,0.752) (-4.046,0.799) (-4.052,0.847) (-4.058,0.895) (-4.064,0.942) (-4.070,0.990) (-4.077,1.037) (-4.083,1.085) (-4.090,1.132) (-4.097,1.180) (-4.104,1.227) (-4.111,1.275) (-4.118,1.322) (-4.126,1.370) (-4.134,1.417) (-4.142,1.464) (-4.150,1.512) (-4.158,1.559) (-4.166,1.606) (-4.175,1.654) (-4.184,1.701) (-4.193,1.748) (-4.202,1.795) (-4.212,1.842) (-4.221,1.889) (-4.231,1.936) (-4.241,1.983) (-4.251,2.030) (-4.262,2.077) (-4.272,2.124) (-4.283,2.170) (-4.294,2.217) (-4.305,2.264) (-4.316,2.310) (-4.328,2.357) (-4.340,2.404) (-4.351,2.450) (-4.363,2.497) (-4.376,2.543) (-4.388,2.589) (-4.401,2.636) (-4.413,2.682) (-4.426,2.728) (-4.439,2.774) (-4.453,2.821) (-4.466,2.867) (-4.480,2.913) (-4.490,2.947)};
\draw[flinea, flecha=1.166cm] plot coordinates {(-3.968,-0.084) (-3.973,-0.131) (-3.978,-0.179) (-3.983,-0.227) (-3.988,-0.275) (-3.993,-0.322) (-3.998,-0.370) (-4.003,-0.418) (-4.008,-0.466) (-4.013,-0.513) (-4.018,-0.561) (-4.024,-0.609) (-4.029,-0.656) (-4.035,-0.704) (-4.040,-0.752) (-4.046,-0.799) (-4.052,-0.847) (-4.058,-0.895) (-4.064,-0.942) (-4.070,-0.990) (-4.077,-1.037) (-4.083,-1.085) (-4.090,-1.132) (-4.097,-1.180) (-4.104,-1.227) (-4.111,-1.275) (-4.118,-1.322) (-4.126,-1.370) (-4.134,-1.417) (-4.142,-1.464) (-4.150,-1.512) (-4.158,-1.559) (-4.166,-1.606) (-4.175,-1.654) (-4.184,-1.701) (-4.193,-1.748) (-4.202,-1.795) (-4.212,-1.842) (-4.221,-1.889) (-4.231,-1.936) (-4.241,-1.983) (-4.251,-2.030) (-4.262,-2.077) (-4.272,-2.124) (-4.283,-2.170) (-4.294,-2.217) (-4.305,-2.264) (-4.316,-2.310) (-4.328,-2.357) (-4.340,-2.404) (-4.351,-2.450) (-4.363,-2.497) (-4.376,-2.543) (-4.388,-2.589) (-4.401,-2.636) (-4.413,-2.682) (-4.426,-2.728) (-4.439,-2.774) (-4.453,-2.821) (-4.466,-2.867) (-4.480,-2.913) (-4.490,-2.947)};
\draw[flinea, flecha=1.234cm] plot coordinates {(-3.985,0.080) (-4.000,0.126) (-4.014,0.172) (-4.029,0.217) (-4.043,0.263) (-4.058,0.309) (-4.072,0.355) (-4.087,0.400) (-4.101,0.446) (-4.116,0.492) (-4.131,0.537) (-4.146,0.583) (-4.161,0.629) (-4.176,0.674) (-4.191,0.720) (-4.206,0.765) (-4.221,0.811) (-4.237,0.856) (-4.252,0.902) (-4.268,0.947) (-4.284,0.992) (-4.300,1.038) (-4.316,1.083) (-4.332,1.128) (-4.348,1.173) (-4.364,1.219) (-4.381,1.264) (-4.398,1.309) (-4.414,1.354) (-4.431,1.399) (-4.448,1.443) (-4.465,1.488) (-4.483,1.533) (-4.500,1.578) (-4.518,1.622) (-4.536,1.667) (-4.554,1.711) (-4.572,1.756) (-4.590,1.800) (-4.608,1.845) (-4.626,1.889) (-4.645,1.933) (-4.664,1.978) (-4.683,2.022) (-4.702,2.066) (-4.721,2.110) (-4.740,2.154) (-4.759,2.198) (-4.779,2.242) (-4.799,2.285) (-4.818,2.329) (-4.838,2.373) (-4.858,2.416) (-4.878,2.460) (-4.899,2.503) (-4.919,2.547) (-4.940,2.590) (-4.960,2.634) (-4.981,2.677) (-5.002,2.720) (-5.023,2.763) (-5.044,2.806) (-5.065,2.850) (-5.086,2.893) (-5.108,2.936) (-5.113,2.946)};
\draw[flinea, flecha=1.234cm] plot coordinates {(-3.985,-0.080) (-4.000,-0.126) (-4.014,-0.172) (-4.029,-0.217) (-4.043,-0.263) (-4.058,-0.309) (-4.072,-0.355) (-4.087,-0.400) (-4.101,-0.446) (-4.116,-0.492) (-4.131,-0.537) (-4.146,-0.583) (-4.161,-0.629) (-4.176,-0.674) (-4.191,-0.720) (-4.206,-0.765) (-4.221,-0.811) (-4.237,-0.856) (-4.252,-0.902) (-4.268,-0.947) (-4.284,-0.992) (-4.300,-1.038) (-4.316,-1.083) (-4.332,-1.128) (-4.348,-1.173) (-4.364,-1.219) (-4.381,-1.264) (-4.398,-1.309) (-4.414,-1.354) (-4.431,-1.399) (-4.448,-1.443) (-4.465,-1.488) (-4.483,-1.533) (-4.500,-1.578) (-4.518,-1.622) (-4.536,-1.667) (-4.554,-1.711) (-4.572,-1.756) (-4.590,-1.800) (-4.608,-1.845) (-4.626,-1.889) (-4.645,-1.933) (-4.664,-1.978) (-4.683,-2.022) (-4.702,-2.066) (-4.721,-2.110) (-4.740,-2.154) (-4.759,-2.198) (-4.779,-2.242) (-4.799,-2.285) (-4.818,-2.329) (-4.838,-2.373) (-4.858,-2.416) (-4.878,-2.460) (-4.899,-2.503) (-4.919,-2.547) (-4.940,-2.590) (-4.960,-2.634) (-4.981,-2.677) (-5.002,-2.720) (-5.023,-2.763) (-5.044,-2.806) (-5.065,-2.850) (-5.086,-2.893) (-5.108,-2.936) (-5.113,-2.946)};
\draw[flinea, flecha=0.878cm] plot coordinates {(-4.002,0.073) (-4.026,0.114) (-4.050,0.156) (-4.074,0.197) (-4.098,0.239) (-4.122,0.280) (-4.146,0.322) (-4.171,0.363) (-4.195,0.405) (-4.219,0.446) (-4.243,0.488) (-4.268,0.529) (-4.292,0.570) (-4.317,0.612) (-4.341,0.653) (-4.366,0.694) (-4.391,0.735) (-4.415,0.776) (-4.440,0.817) (-4.465,0.858) (-4.490,0.899) (-4.515,0.940) (-4.541,0.981) (-4.566,1.022) (-4.591,1.063) (-4.617,1.103) (-4.642,1.144) (-4.668,1.184) (-4.694,1.225) (-4.719,1.265) (-4.745,1.306) (-4.771,1.346) (-4.798,1.386) (-4.824,1.427) (-4.850,1.467) (-4.876,1.507) (-4.903,1.547) (-4.930,1.587) (-4.956,1.627) (-4.983,1.667) (-5.010,1.706) (-5.037,1.746) (-5.064,1.786) (-5.091,1.825) (-5.118,1.865) (-5.146,1.904) (-5.166,1.934)};
\draw[flinea, flecha=0.878cm] plot coordinates {(-4.002,-0.073) (-4.026,-0.114) (-4.050,-0.156) (-4.074,-0.197) (-4.098,-0.239) (-4.122,-0.280) (-4.146,-0.322) (-4.171,-0.363) (-4.195,-0.405) (-4.219,-0.446) (-4.243,-0.488) (-4.268,-0.529) (-4.292,-0.570) (-4.317,-0.612) (-4.341,-0.653) (-4.366,-0.694) (-4.391,-0.735) (-4.415,-0.776) (-4.440,-0.817) (-4.465,-0.858) (-4.490,-0.899) (-4.515,-0.940) (-4.541,-0.981) (-4.566,-1.022) (-4.591,-1.063) (-4.617,-1.103) (-4.642,-1.144) (-4.668,-1.184) (-4.694,-1.225) (-4.719,-1.265) (-4.745,-1.306) (-4.771,-1.346) (-4.798,-1.386) (-4.824,-1.427) (-4.850,-1.467) (-4.876,-1.507) (-4.903,-1.547) (-4.930,-1.587) (-4.956,-1.627) (-4.983,-1.667) (-5.010,-1.706) (-5.037,-1.746) (-5.064,-1.786) (-5.091,-1.825) (-5.118,-1.865) (-5.146,-1.904) (-5.166,-1.934)};
\draw[flinea, flecha=0.643cm] plot coordinates {(-4.019,0.060) (-4.052,0.094) (-4.086,0.129) (-4.120,0.163) (-4.153,0.197) (-4.187,0.231) (-4.221,0.265) (-4.254,0.300) (-4.288,0.334) (-4.322,0.368) (-4.356,0.402) (-4.390,0.436) (-4.423,0.470) (-4.457,0.504) (-4.491,0.538) (-4.525,0.572) (-4.559,0.605) (-4.593,0.639) (-4.628,0.673) (-4.662,0.707) (-4.696,0.740) (-4.730,0.774) (-4.765,0.807) (-4.799,0.841) (-4.834,0.874) (-4.868,0.908) (-4.903,0.941) (-4.937,0.974) (-4.972,1.007) (-5.007,1.041) (-5.041,1.074) (-5.076,1.107) (-5.111,1.140) (-5.146,1.173) (-5.164,1.189)};
\draw[flinea, flecha=0.643cm] plot coordinates {(-4.019,-0.060) (-4.052,-0.094) (-4.086,-0.129) (-4.120,-0.163) (-4.153,-0.197) (-4.187,-0.231) (-4.221,-0.265) (-4.254,-0.300) (-4.288,-0.334) (-4.322,-0.368) (-4.356,-0.402) (-4.390,-0.436) (-4.423,-0.470) (-4.457,-0.504) (-4.491,-0.538) (-4.525,-0.572) (-4.559,-0.605) (-4.593,-0.639) (-4.628,-0.673) (-4.662,-0.707) (-4.696,-0.740) (-4.730,-0.774) (-4.765,-0.807) (-4.799,-0.841) (-4.834,-0.874) (-4.868,-0.908) (-4.903,-0.941) (-4.937,-0.974) (-4.972,-1.007) (-5.007,-1.041) (-5.041,-1.074) (-5.076,-1.107) (-5.111,-1.140) (-5.146,-1.173) (-5.164,-1.189)};
\draw[flinea, flecha=0.499cm] plot coordinates {(-4.036,0.037) (-4.079,0.058) (-4.122,0.078) (-4.165,0.099) (-4.208,0.120) (-4.252,0.141) (-4.295,0.162) (-4.338,0.183) (-4.381,0.204) (-4.425,0.224) (-4.468,0.245) (-4.511,0.266) (-4.555,0.287) (-4.598,0.307) (-4.641,0.328) (-4.685,0.349) (-4.728,0.369) (-4.771,0.390) (-4.815,0.410) (-4.858,0.431) (-4.901,0.451) (-4.945,0.472) (-4.988,0.492) (-5.032,0.512) (-5.075,0.533) (-5.119,0.553) (-5.162,0.573)};
\draw[flinea, flecha=0.499cm] plot coordinates {(-4.036,-0.037) (-4.079,-0.058) (-4.122,-0.078) (-4.165,-0.099) (-4.208,-0.120) (-4.252,-0.141) (-4.295,-0.162) (-4.338,-0.183) (-4.381,-0.204) (-4.425,-0.224) (-4.468,-0.245) (-4.511,-0.266) (-4.555,-0.287) (-4.598,-0.307) (-4.641,-0.328) (-4.685,-0.349) (-4.728,-0.369) (-4.771,-0.390) (-4.815,-0.410) (-4.858,-0.431) (-4.901,-0.451) (-4.945,-0.472) (-4.988,-0.492) (-5.032,-0.512) (-5.075,-0.533) (-5.119,-0.553) (-5.162,-0.573)};
\draw[flinea, flecha=0.945cm] plot coordinates {(1.625,0.787) (1.668,0.808) (1.712,0.829) (1.755,0.850) (1.798,0.871) (1.842,0.892) (1.885,0.912) (1.928,0.933) (1.972,0.953) (2.015,0.974) (2.058,0.994) (2.102,1.015) (2.145,1.035) (2.189,1.055) (2.232,1.076) (2.276,1.096) (2.319,1.116) (2.363,1.136) (2.407,1.156) (2.450,1.176) (2.494,1.196) (2.538,1.216) (2.581,1.236) (2.625,1.256) (2.669,1.276) (2.712,1.295) (2.756,1.315) (2.800,1.335) (2.844,1.354) (2.888,1.374) (2.932,1.393) (2.975,1.413) (3.019,1.432) (3.063,1.452) (3.107,1.471) (3.151,1.490) (3.195,1.510) (3.239,1.529) (3.283,1.548) (3.327,1.567) (3.371,1.586) (3.415,1.605) (3.459,1.624) (3.503,1.643) (3.547,1.662) (3.591,1.681) (3.636,1.700) (3.680,1.719) (3.724,1.738) (3.768,1.757) (3.812,1.775) (3.856,1.794) (3.901,1.813) (3.945,1.832) (3.989,1.850) (4.033,1.869) (4.078,1.887) (4.089,1.892)};
\draw[flinea, flecha=0.945cm] plot coordinates {(1.625,-0.787) (1.668,-0.808) (1.712,-0.829) (1.755,-0.850) (1.798,-0.871) (1.842,-0.892) (1.885,-0.912) (1.928,-0.933) (1.972,-0.953) (2.015,-0.974) (2.058,-0.994) (2.102,-1.015) (2.145,-1.035) (2.189,-1.055) (2.232,-1.076) (2.276,-1.096) (2.319,-1.116) (2.363,-1.136) (2.407,-1.156) (2.450,-1.176) (2.494,-1.196) (2.538,-1.216) (2.581,-1.236) (2.625,-1.256) (2.669,-1.276) (2.712,-1.295) (2.756,-1.315) (2.800,-1.335) (2.844,-1.354) (2.888,-1.374) (2.932,-1.393) (2.975,-1.413) (3.019,-1.432) (3.063,-1.452) (3.107,-1.471) (3.151,-1.490) (3.195,-1.510) (3.239,-1.529) (3.283,-1.548) (3.327,-1.567) (3.371,-1.586) (3.415,-1.605) (3.459,-1.624) (3.503,-1.643) (3.547,-1.662) (3.591,-1.681) (3.636,-1.700) (3.680,-1.719) (3.724,-1.738) (3.768,-1.757) (3.812,-1.775) (3.856,-1.794) (3.901,-1.813) (3.945,-1.832) (3.989,-1.850) (4.033,-1.869) (4.078,-1.887) (4.089,-1.892)};
\draw[flinea, flecha=0.853cm] plot coordinates {(1.259,1.295) (1.292,1.329) (1.326,1.364) (1.359,1.398) (1.393,1.432) (1.427,1.466) (1.461,1.500) (1.495,1.534) (1.529,1.568) (1.563,1.602) (1.597,1.635) (1.632,1.669) (1.666,1.702) (1.701,1.736) (1.735,1.769) (1.770,1.802) (1.804,1.835) (1.839,1.869) (1.874,1.902) (1.909,1.934) (1.944,1.967) (1.979,2.000) (2.014,2.033) (2.049,2.065) (2.085,2.098) (2.120,2.130) (2.156,2.163) (2.191,2.195) (2.227,2.227) (2.262,2.259) (2.298,2.291) (2.334,2.323) (2.370,2.355) (2.405,2.387) (2.441,2.419) (2.477,2.451) (2.513,2.483) (2.550,2.514) (2.586,2.546) (2.622,2.577) (2.658,2.609) (2.695,2.640) (2.731,2.671) (2.767,2.703) (2.804,2.734) (2.840,2.765) (2.877,2.796) (2.914,2.827) (2.950,2.858) (2.987,2.889) (3.024,2.920) (3.051,2.943)};
\draw[flinea, flecha=0.853cm] plot coordinates {(1.259,-1.295) (1.292,-1.329) (1.326,-1.364) (1.359,-1.398) (1.393,-1.432) (1.427,-1.466) (1.461,-1.500) (1.495,-1.534) (1.529,-1.568) (1.563,-1.602) (1.597,-1.635) (1.632,-1.669) (1.666,-1.702) (1.701,-1.736) (1.735,-1.769) (1.770,-1.802) (1.804,-1.835) (1.839,-1.869) (1.874,-1.902) (1.909,-1.934) (1.944,-1.967) (1.979,-2.000) (2.014,-2.033) (2.049,-2.065) (2.085,-2.098) (2.120,-2.130) (2.156,-2.163) (2.191,-2.195) (2.227,-2.227) (2.262,-2.259) (2.298,-2.291) (2.334,-2.323) (2.370,-2.355) (2.405,-2.387) (2.441,-2.419) (2.477,-2.451) (2.513,-2.483) (2.550,-2.514) (2.586,-2.546) (2.622,-2.577) (2.658,-2.609) (2.695,-2.640) (2.731,-2.671) (2.767,-2.703) (2.804,-2.734) (2.840,-2.765) (2.877,-2.796) (2.914,-2.827) (2.950,-2.858) (2.987,-2.889) (3.024,-2.920) (3.051,-2.943)};
\draw[flinea, flecha=0.571cm] plot coordinates {(0.883,1.575) (0.907,1.617) (0.931,1.659) (0.955,1.700) (0.979,1.742) (1.003,1.783) (1.027,1.825) (1.052,1.866) (1.076,1.907) (1.101,1.948) (1.126,1.989) (1.151,2.030) (1.177,2.071) (1.202,2.112) (1.228,2.152) (1.254,2.193) (1.280,2.233) (1.306,2.273) (1.332,2.313) (1.358,2.354) (1.385,2.394) (1.411,2.433) (1.438,2.473) (1.465,2.513) (1.492,2.553) (1.519,2.592) (1.547,2.632) (1.574,2.671) (1.602,2.710) (1.630,2.750) (1.657,2.789) (1.685,2.828) (1.713,2.867) (1.741,2.906) (1.770,2.944)};
\draw[flinea, flecha=0.571cm] plot coordinates {(0.883,-1.575) (0.907,-1.617) (0.931,-1.659) (0.955,-1.700) (0.979,-1.742) (1.003,-1.783) (1.027,-1.825) (1.052,-1.866) (1.076,-1.907) (1.101,-1.948) (1.126,-1.989) (1.151,-2.030) (1.177,-2.071) (1.202,-2.112) (1.228,-2.152) (1.254,-2.193) (1.280,-2.233) (1.306,-2.273) (1.332,-2.313) (1.358,-2.354) (1.385,-2.394) (1.411,-2.433) (1.438,-2.473) (1.465,-2.513) (1.492,-2.553) (1.519,-2.592) (1.547,-2.632) (1.574,-2.671) (1.602,-2.710) (1.630,-2.750) (1.657,-2.789) (1.685,-2.828) (1.713,-2.867) (1.741,-2.906) (1.770,-2.944)};
\draw[flinea, flecha=0.454cm] plot coordinates {(0.495,1.737) (0.509,1.783) (0.522,1.829) (0.536,1.875) (0.550,1.921) (0.565,1.966) (0.579,2.012) (0.594,2.058) (0.610,2.103) (0.625,2.149) (0.641,2.194) (0.657,2.239) (0.673,2.284) (0.690,2.329) (0.707,2.374) (0.724,2.419) (0.741,2.464) (0.759,2.509) (0.777,2.553) (0.795,2.598) (0.813,2.642) (0.832,2.686) (0.850,2.731) (0.869,2.775) (0.888,2.819) (0.908,2.863) (0.927,2.906) (0.947,2.950)};
\draw[flinea, flecha=0.454cm] plot coordinates {(0.495,-1.737) (0.509,-1.783) (0.522,-1.829) (0.536,-1.875) (0.550,-1.921) (0.565,-1.966) (0.579,-2.012) (0.594,-2.058) (0.610,-2.103) (0.625,-2.149) (0.641,-2.194) (0.657,-2.239) (0.673,-2.284) (0.690,-2.329) (0.707,-2.374) (0.724,-2.419) (0.741,-2.464) (0.759,-2.509) (0.777,-2.553) (0.795,-2.598) (0.813,-2.642) (0.832,-2.686) (0.850,-2.731) (0.869,-2.775) (0.888,-2.819) (0.908,-2.863) (0.927,-2.906) (0.947,-2.950)};
\draw[flinea, flecha=0.407cm] plot coordinates {(0.088,1.804) (0.090,1.852) (0.093,1.900) (0.097,1.948) (0.101,1.995) (0.105,2.043) (0.110,2.091) (0.115,2.139) (0.121,2.186) (0.128,2.234) (0.134,2.281) (0.141,2.329) (0.149,2.376) (0.157,2.424) (0.165,2.471) (0.174,2.518) (0.183,2.565) (0.193,2.612) (0.203,2.659) (0.213,2.706) (0.223,2.753) (0.234,2.800) (0.246,2.846) (0.257,2.893) (0.269,2.939) (0.272,2.951)};
\draw[flinea, flecha=0.407cm] plot coordinates {(0.088,-1.804) (0.090,-1.852) (0.093,-1.900) (0.097,-1.948) (0.101,-1.995) (0.105,-2.043) (0.110,-2.091) (0.115,-2.139) (0.121,-2.186) (0.128,-2.234) (0.134,-2.281) (0.141,-2.329) (0.149,-2.376) (0.157,-2.424) (0.165,-2.471) (0.174,-2.518) (0.183,-2.565) (0.193,-2.612) (0.203,-2.659) (0.213,-2.706) (0.223,-2.753) (0.234,-2.800) (0.246,-2.846) (0.257,-2.893) (0.269,-2.939) (0.272,-2.951)};
\draw[flinea, flecha=0.412cm] plot coordinates {(-0.356,1.771) (-0.365,1.818) (-0.373,1.865) (-0.380,1.912) (-0.387,1.960) (-0.393,2.008) (-0.398,2.055) (-0.402,2.103) (-0.405,2.151) (-0.408,2.199) (-0.410,2.247) (-0.411,2.295) (-0.412,2.343) (-0.411,2.391) (-0.411,2.439) (-0.409,2.487) (-0.407,2.535) (-0.405,2.583) (-0.402,2.631) (-0.398,2.679) (-0.394,2.726) (-0.390,2.774) (-0.385,2.822) (-0.380,2.870) (-0.374,2.917) (-0.371,2.941)};
\draw[flinea, flecha=0.412cm] plot coordinates {(-0.356,-1.771) (-0.365,-1.818) (-0.373,-1.865) (-0.380,-1.912) (-0.387,-1.960) (-0.393,-2.008) (-0.398,-2.055) (-0.402,-2.103) (-0.405,-2.151) (-0.408,-2.199) (-0.410,-2.247) (-0.411,-2.295) (-0.412,-2.343) (-0.411,-2.391) (-0.411,-2.439) (-0.409,-2.487) (-0.407,-2.535) (-0.405,-2.583) (-0.402,-2.631) (-0.398,-2.679) (-0.394,-2.726) (-0.390,-2.774) (-0.385,-2.822) (-0.380,-2.870) (-0.374,-2.917) (-0.371,-2.941)};
\draw[flinea, flecha=0.491cm] plot coordinates {(-0.892,1.570) (-0.915,1.612) (-0.936,1.656) (-0.955,1.700) (-0.972,1.745) (-0.987,1.790) (-1.000,1.836) (-1.012,1.883) (-1.022,1.930) (-1.031,1.977) (-1.039,2.024) (-1.046,2.072) (-1.051,2.120) (-1.056,2.167) (-1.059,2.215) (-1.062,2.263) (-1.064,2.311) (-1.065,2.359) (-1.066,2.407) (-1.066,2.455) (-1.066,2.503) (-1.065,2.551) (-1.063,2.599) (-1.062,2.647) (-1.059,2.695) (-1.057,2.743) (-1.054,2.791) (-1.050,2.839) (-1.047,2.886) (-1.043,2.934) (-1.042,2.946)};
\draw[flinea, flecha=0.491cm] plot coordinates {(-0.892,-1.570) (-0.915,-1.612) (-0.936,-1.656) (-0.955,-1.700) (-0.972,-1.745) (-0.987,-1.790) (-1.000,-1.836) (-1.012,-1.883) (-1.022,-1.930) (-1.031,-1.977) (-1.039,-2.024) (-1.046,-2.072) (-1.051,-2.120) (-1.056,-2.167) (-1.059,-2.215) (-1.062,-2.263) (-1.064,-2.311) (-1.065,-2.359) (-1.066,-2.407) (-1.066,-2.455) (-1.066,-2.503) (-1.065,-2.551) (-1.063,-2.599) (-1.062,-2.647) (-1.059,-2.695) (-1.057,-2.743) (-1.054,-2.791) (-1.050,-2.839) (-1.047,-2.886) (-1.043,-2.934) (-1.042,-2.946)};

    % Generado automáticamente por generar_figuras.py; no editar.
\node[menos] at (1.276,0.384) {};
\node[menos] at (1.276,-0.384) {};
\node[menos] at (1.139,0.690) {};
\node[menos] at (1.139,-0.690) {};
\node[menos] at (0.961,0.922) {};
\node[menos] at (0.961,-0.922) {};
\node[menos] at (0.722,1.119) {};
\node[menos] at (0.722,-1.119) {};
\node[menos] at (0.390,1.274) {};
\node[menos] at (0.390,-1.274) {};
\node[menos] at (-0.089,1.329) {};
\node[menos] at (-0.089,-1.329) {};
\node[menos] at (-0.817,1.052) {};
\node[menos] at (-0.817,-1.052) {};
\node[mas] at (1.501,0.727) {};
\node[mas] at (1.501,-0.727) {};
\node[mas] at (1.163,1.196) {};
\node[mas] at (1.163,-1.196) {};
\node[mas] at (0.816,1.455) {};
\node[mas] at (0.816,-1.455) {};
\node[mas] at (0.458,1.604) {};
\node[mas] at (0.458,-1.604) {};
\node[mas] at (0.081,1.666) {};
\node[mas] at (0.081,-1.666) {};
\node[mas] at (-0.329,1.635) {};
\node[mas] at (-0.329,-1.635) {};
\node[mas] at (-0.824,1.450) {};
\node[mas] at (-0.824,-1.450) {};

    \node[cargapos] at (0.38*\esc,0) {};
    \node[cargapos] at (-3.3*\esc,0) {};
    \node[rotulo,fill=fondo,inner sep=1pt,rounded corners=1pt,anchor=south] at (0.38*\esc,0.17) {$q$};
    \node[rotulo,fill=panel,inner sep=1pt,rounded corners=1pt,anchor=south] at (-3.3*\esc,0.18) {$q_{\mathrm{e}}$};
  \end{scope}
  \draw[azulborde,line width=0.5pt] (-4.3*\esc,-2.45*\esc) rectangle (3.4*\esc,2.45*\esc);
  \node[rotulogris,anchor=south west] at (-4.3*\esc,2.47*\esc)
     {(a) Corteza neutra: carga en la cavidad y carga exterior};
  % ---------------- (b)
  \begin{scope}[shift={(6.6,0)}]
    \fill[panel] (-1.85,-2.94) rectangle (2.35,2.94);
    \fill[even odd rule,inner color=conductor,outer color=conductorosc]
      (0,0) circle (1.6) (0,0) circle (0.95);
    \draw[gristexto,line width=0.5pt] (0,0) circle (1.6) (0,0) circle (0.95);
    \fill[fondo] (0,0) circle (0.95);
    \draw[vectorgris,line width=0.45pt] (0,0) -- (135:0.95) node[pos=0.55,above right=-1pt,rotulo] {$a$};
    \draw[vectorgris,line width=0.45pt] (0,0) -- (-35:1.6) node[pos=0.78,above right=-1pt,rotulo] {$b$};
    \node[punto] at (0,0) {};
    \node[rotulo,text=azul] at (-0.3,-0.45) {$V$};
    % conexión a tierra
    \draw[gristexto,line width=0.6pt] (1.6,0) -- (2.0,0) -- (2.0,-0.55);
    \tierra{(2.0,-0.55)}
    \node[rotulo,align=center,text=oro] at (0.15,-2.35) {$\varphi\big|_{r=a}=0\ \Rightarrow\ \sigma_a=0$};
    \draw[azulborde,line width=0.5pt] (-1.85,-2.94) rectangle (2.35,2.94);
    \node[rotulogris,anchor=south west] at (-1.85,2.96) {(b) Cavidad vacía};
  \end{scope}
\end{tikzpicture}
